\chapter{Implementation}
\label{ch:implementation}

\section{The from-scratch constraint}

Both algorithms and the environment are written from first principles in NumPy. No reinforcement
learning library is used: there is no \code{gymnasium}, no \code{stable-baselines}, and no
pre-built agent. NumPy supplies array arithmetic and random number generation; everything
specific to the MDP, to temporal-difference control, and to the two update rules is written out.

\section{The environment}

\code{WarehouseEnv} is a plain class with \code{reset()} and \code{step()}.

{\small
\begin{verbatim}
def step(self, action):
    self.t += 1
    dr, dc = MOVES[action]
    nr, nc = self.pos[0] + dr, self.pos[1] + dc
    info = {}

    if not (0 <= nr < self.H and 0 <= nc < self.W) or self.grid[nr, nc] == 1:
        reward, done = self.r_bump, False        # nudge: no movement
        info["bump"] = True
    elif self.walkway[nr, nc]:
        self.pos = self.start                    # incident: back to charger
        reward, done = self.r_incident, False
        info["incident"] = True
    else:
        self.pos = (nr, nc)
        if self.pos == self.goal:
            reward, done, info["goal"] = self.r_goal, True, True
        else:
            reward, done = self.r_step, False
    if self.t >= self.max_steps:
        done, info["timeout"] = True, True
    return self.encode(self.pos), reward, done, info
\end{verbatim}}

The three outcomes are mutually exclusive and ordered deliberately: a cell that is both racking and
walkway cannot occur by construction, and the boundary check precedes the grid lookup so that an
out-of-bounds index is never formed.

State encoding is $s = r \cdot W + c$ with the inverse $(s \div W,\; s \bmod W)$. Bijectivity over
all 400 cells is asserted in Chapter~\ref{ch:verification}.

\section{The shared agent}

The two algorithms share a base class holding everything they have in common: the action-value
table, $\varepsilon$-greedy action selection with random tie-breaking, and the $\varepsilon$
schedule.

{\small
\begin{verbatim}
class TDControlAgent:
    def act(self, s):
        if self.rng.random() < self.eps:
            return int(self.rng.integers(self.n_actions))
        return self.greedy(s)

    def greedy(self, s):
        q = self.Q[s]
        return int(self.rng.choice(np.flatnonzero(q == q.max())))   # random tie-break

    def update(self, s, a, r, s_next, a_next, done):
        tgt = self._target(r, s_next, a_next, done)
        self.Q[s, a] += self.alpha * (tgt - self.Q[s, a])
\end{verbatim}}

Random tie-breaking matters more than it appears. With $Q$ initialised to zeros every action ties
initially, and breaking ties by \code{argmax} would bias the agent toward action $0$ --- ``up'',
directly into the walkway --- producing an artefact that looks like a finding.

\section{The one-line difference}

The subclasses override a single method. This is the entire distinction between the two
algorithms, and the code is arranged so that nothing else can differ:

{\small
\begin{verbatim}
class QLearningAgent(TDControlAgent):
    # Off-policy: bootstraps from the best available next action.
    name = "Q-learning"
    def _target(self, r, s_next, a_next, done):
        return r if done else r + self.gamma * self.Q[s_next].max()


class SarsaAgent(TDControlAgent):
    # On-policy: bootstraps from the action actually selected for the next step.
    name = "SARSA"
    def _target(self, r, s_next, a_next, done):
        return r if done else r + self.gamma * self.Q[s_next, a_next]
\end{verbatim}}

These implement Equations~\ref{eq:qlearning} and~\ref{eq:sarsa} respectively. Note the terminal
handling: when \code{done} is true the target is $r$ alone, with no bootstrap term, since a
terminal state has value zero by definition. Omitting that condition is a common error which
inflates all values by a spurious $\gamma V(s_{\text{terminal}})$ term.

\section{The episode loop}

{\small
\begin{verbatim}
def run_episode(env, agent, learn=True, greedy=False):
    s = env.reset()
    a = agent.greedy(s) if greedy else agent.act(s)
    while True:
        s_next, r, done, info = env.step(a)
        a_next = (agent.greedy(s_next) if greedy
                  else agent.act(s_next)) if not done else 0
        if learn:
            agent.update(s, a, r, s_next, a_next, done)
        if done: break
        s, a = s_next, a_next
\end{verbatim}}

\keypoint{The critical detail is that \code{a\us next} is chosen before the update and then
actually executed} on the following iteration, via \code{s, a = s\us next, a\us next}. SARSA is
on-policy precisely because the action it bootstraps from is the action it takes. Sampling a fresh
action after the update --- an easy and common slip --- silently converts SARSA into an
inconsistent hybrid that still trains and still produces a plausible learning curve. The
verification suite is designed to catch exactly this.

\section{Experimental harness}

A single \code{train()} function serves both algorithms, so the training budget is identical by
construction rather than by careful duplication.

\begin{table}[ht]
\centering
\caption[Hyperparameters]{Hyperparameters. Every value is shared between the two algorithms.}
\label{tab:hyper}
\begin{tabular}{lll}
\toprule
\textbf{Parameter} & \textbf{Value} & \textbf{Note} \\
\midrule
Learning rate $\alpha$      & $0.10$ & constant \\
Discount $\gamma$           & $0.99$ & \\
$\varepsilon$ schedule      & $1.0 \rightarrow 0.05$ & linear over the first 60\% of training \\
Episodes per run            & 3{,}000 & \\
Episode cap                 & 600 steps & \\
Seeds                       & 30 & paired across algorithms \\
\bottomrule
\end{tabular}
\end{table}

The residual $\varepsilon = 0.05$ is deliberate. Driving $\varepsilon$ to zero would make SARSA
converge to the same policy as Q-learning, eliminating the effect under study by construction.

Each seed trains both a Q-learning and a SARSA agent on the same floor with the same budget, the
same hyperparameters and the same schedule. Evaluation is then performed separately with
$\varepsilon = 0$, which distinguishes the policy the agent \emph{learned} from the behaviour it
exhibited while learning --- the distinction on which the entire study turns.
